\chapter{Pecahan Pertama}\label{ch:g4:fractions}

Setengah buah apel, \omterm{def:g3:division:half}{seperempat} jam, dan tiga perempat isi kelas:
\omterm{def:g4:fractions:def}{pecahan}
menamai potongan sesuatu. Bab ini memperkenalkan penulisan
$\frac{a}{b}$, membaca \omterm{def:g4:fractions:def}{pecahan} pada gambar dan garis bilangan, lalu
membandingkan \omterm{def:g4:fractions:def}{pecahan} yang paling sederhana. (\omterm{def:g4:fractions:def}{Pecahan} tumbuh dewasa
pada \cref{ch:g6:fractions}.)

\section{Menamai potongannya}

\begin{definition}[Pecahan]\label{def:g4:fractions:def}
Potonglah sebuah keseluruhan menjadi bagian yang sama besar lalu
ambil beberapa --- itulah \emph{pecahan}\index{pecahan}:
\[
\frac{3}{4}
\quad
\begin{array}{l}
3 \text{: berapa bagian yang diambil (\emph{pembilang}\index{pembilang});}\\
4 \text{: keseluruhan dipotong menjadi berapa bagian yang sama besar
(\emph{penyebut}\index{penyebut}).}
\end{array}
\]
$\frac12$ adalah \emph{setengah}, $\frac13$ \emph{sepertiga},
$\frac14$ \emph{\omterm{def:g3:division:half}{seperempat}}, $\frac{1}{10}$ \emph{sepersepuluh}.
\end{definition}

\begin{omfigure}
\begin{tikzpicture}[scale=0.78]
  % setengah lingkaran
  \draw[thick] (0,0) circle (1.0);
  \fill[omDef!40] (0,0) -- (1,0) arc (0:180:1) -- cycle;
  \draw[thick] (-1,0) -- (1,0);
  \node[below, font=\small] at (0,-1.2) {$\frac12$};
  % seperempat
  \begin{scope}[xshift=3.6cm]
  \draw[thick] (0,0) circle (1.0);
  \fill[omThm!40] (0,0) -- (1,0) arc (0:90:1) -- cycle;
  \fill[omThm!40] (0,0) -- (0,1) arc (90:180:1) -- cycle;
  \fill[omThm!40] (0,0) -- (-1,0) arc (180:270:1) -- cycle;
  \draw[thick] (-1,0) -- (1,0);
  \draw[thick] (0,-1) -- (0,1);
  \node[below, font=\small] at (0,-1.2) {$\frac34$};
  \end{scope}
  % batang sepertiga
  \begin{scope}[xshift=6.6cm, yshift=-0.35cm]
  \foreach \i in {0,1,2} \draw[omExo] (\i*1.3,0) rectangle
    (\i*1.3+1.3,0.7);
  \fill[omMeth!40] (0.04,0.04) rectangle (1.26,0.66);
  \draw[thick] (0,0) rectangle (3.9,0.7);
  \node[below, font=\small] at (1.95,-0.15) {$\frac13$};
  \end{scope}
  % batang sepersepuluh
  \begin{scope}[xshift=11.6cm, yshift=-0.35cm]
  \foreach \i in {0,...,9} \draw[omExo] (\i*0.42,0) rectangle
    (\i*0.42+0.42,0.7);
  \foreach \i in {0,...,6} \fill[omProp!40] (\i*0.42+0.03,0.04)
    rectangle (\i*0.42+0.39,0.66);
  \draw[thick] (0,0) rectangle (4.2,0.7);
  \node[below, font=\small] at (2.1,-0.15) {$\frac{7}{10}$};
  \end{scope}
\end{tikzpicture}

{\small Membaca \omterm{def:g4:fractions:def}{pecahan} pada gambar. Bagiannya harus \emph{sama
besar}: tiga potong yang tidak sama besar bukanlah sepertiga!}
\end{omfigure}

\begin{example}[Pecahan dari sekumpulan benda]\label{ex:g4:fractions:collection}
$\frac13$ dari $12$ kelereng: bagilah $12$ kelereng itu menjadi $3$
kelompok sama besar berisi $4$; satu kelompok adalah $\frac13$, jadi
$\frac13$ dari $12$ adalah $4$ --- dan $\frac23$ dari $12$ adalah
dua kelompok, yaitu $8$.
\end{example}

\section{Pecahan pada garis bilangan}

\begin{method}[Meletakkan $\frac{a}{b}$]\label{met:g4:fractions:line}
Potonglah setiap satuan pada garis menjadi $b$ langkah yang sama;
mulai dari $0$, berjalanlah $a$ langkah. Jika $a$ lebih besar
daripada $b$, jalannya melewati $1$: \omterm{def:g4:fractions:def}{pecahan} dapat lebih besar
daripada satu keseluruhan!
\end{method}

\begin{omfigure}
\begin{tikzpicture}[scale=1.6]
  \draw[-Stealth] (-0.2,0) -- (6.9,0);
  \foreach \i in {0,4,8,12,16,20,24}
    \draw ({\i*0.25},-0.09) -- ({\i*0.25},0.09);
  \foreach \i in {0,...,24} \draw ({\i*0.25},-0.05) -- ({\i*0.25},0.05);
  \foreach \x/\l in {0/0, 1/1, 2/2, 3/3, 4/4, 5/5, 6/6}
    \node[below, font=\small] at (\x,-0.1) {$\l$};
  \fill[omDef] (0.5,0) circle (1.6pt)
    node[above, font=\small] {$\frac24$};
  \fill[omThm] (1.25,0) circle (1.6pt)
    node[above, font=\small] {$\frac54$};
  \fill[omMeth] (2.75,0) circle (1.6pt)
    node[above, font=\small] {$\frac{11}{4}$};
\end{tikzpicture}

{\small Garis yang dipotong menjadi seperempatan: $\frac24$ mendarat
di titik yang sama dengan $\frac12$; $\frac54$ satu \omterm{def:g3:division:half}{seperempat}
melewati $1$; $\frac{11}{4}$ adalah $2$ dan $\frac34$.}
\end{omfigure}

\begin{example}[Bilangan bulat yang bersembunyi di pecahan]\label{ex:g4:fractions:whole}
$\frac44 = 1$ (empat \omterm{def:g3:division:half}{seperempat} menjadi satu keseluruhan);
$\frac82 = 4$; $\frac{12}{3} = 4$. Dan $\frac74$ adalah
$1 + \frac34$: satu keseluruhan ditambah tiga perempat.
\end{example}

\section{Membandingkan dan menjumlahkan pecahan sederhana}

\begin{proposition}[Penyebut yang sama]\label{prop:g4:fractions:compare}
Dengan \omterm{def:g4:fractions:def}{penyebut} yang sama, potongannya sama besar --- jadi cukup
dibilang saja:
\[
\frac{2}{8} < \frac{5}{8},
\qquad
\frac{3}{8} + \frac{4}{8} = \frac{7}{8}.
\]
Membandingkannya dengan satu keseluruhan juga mudah: $\frac{5}{8} < 1 < \frac{9}{8}$
(kurang dari, lalu lebih dari, $8$ perdelapan).
\end{proposition}

\begin{proof}[Mengapa, dalam gambar]
Mewarnai $3$ perdelapan lalu $4$ perdelapan lagi pada batang yang
sama berarti mewarnai $7$ perdelapan seluruhnya.
\end{proof}

\begin{example}[Penyebut berbeda, titik yang sama]\label{ex:g4:fractions:equal}
Pada garis bilangan, $\frac12$, $\frac24$ dan $\frac{5}{10}$
mendarat di titik yang sama: ketiganya adalah tiga nama untuk
bilangan yang sama. Memotong setiap setengah menjadi dua
menghasilkan \omterm{def:g3:division:half}{seperempat}; menjadi lima, sepersepuluh. (Aturan umumnya
ada pada \cref{ch:g6:fractions}.)
\end{example}

\begin{example}[Seperempat jam]\label{ex:g4:fractions:hour}
Satu jam berisi $60$ menit, jadi \omterm{def:g3:division:half}{seperempat} jam adalah
$60 \div 4 = 15$ menit, dan tiga perempat jam adalah $45$ menit.
``Pukul setengah tiga'' itu berarti $\frac12$ jam sesudah pukul dua.
\end{example}

\section{Latihan}

\begin{exercise}[$\star$]\label{exo:g4:fractions:1}
Tulislah \omterm{def:g4:fractions:def}{pecahan} yang ditunjukkan: sebuah pizza dipotong menjadi $6$
dan tersisa $5$ potong; sebatang cokelat berisi $8$ kotak, $3$ sudah
dimakan (\omterm{def:g4:fractions:def}{pecahan} yang dimakan); sebuah bendera dibagi menjadi $3$
jalur tegak yang sama, $1$ diberi warna.
\end{exercise}

\begin{exercise}[$\star$]\label{exo:g4:fractions:2}
Gambarlah sebuah batang yang dipotong menjadi $5$ bagian sama besar
lalu warnailah $\frac35$. Lalu warnai $\frac{2}{5}$ pada batang lain yang sama. Warna mana yang lebih besar?
\end{exercise}

\begin{exercise}[$\star$]\label{exo:g4:fractions:3}
Hitunglah: $\frac12$ dari $18$; \quad $\frac14$ dari $20$; \quad
$\frac13$ dari $21$; \quad $\frac34$ dari $20$ (tiga kelompok
\omterm{def:g3:division:half}{seperempat}).
\end{exercise}

\begin{exercise}[$\star$]\label{exo:g4:fractions:4}
Letakkan pada garis bilangan yang dipotong menjadi sepertigaan:
$\frac13$; \ $\frac33$; \ $\frac53$; \ $2$. \omterm{def:g4:fractions:def}{Pecahan} mana yang mendarat tepat di $2$?
\end{exercise}

\begin{exercise}[$\star$]\label{exo:g4:fractions:5}
Salin dan lengkapi dengan $<$, $>$ atau $=$:
\[
\frac38 \;?\; \frac68, \qquad
\frac55 \;?\; 1, \qquad
\frac74 \;?\; 1, \qquad
\frac12 \;?\; \frac24 .
\]
\end{exercise}

\begin{exercise}[$\star$]\label{exo:g4:fractions:6}
Hitunglah: $\frac26 + \frac36$; \quad $\frac58 - \frac28$; \quad
$\frac14 + \frac24$; \quad $1 - \frac13$ (berapa sepertiga yang
menjadi $1$?).
\end{exercise}

\begin{exercise}[$\star$]\label{exo:g4:fractions:7}
Tulislah sebagai bilangan bulat ditambah \omterm{def:g4:fractions:def}{pecahan} yang lebih kecil
daripada $1$: $\frac54$; \quad $\frac73$; \quad $\frac{9}{2}$.
\end{exercise}

\begin{exercise}[$\star$]\label{exo:g4:fractions:8}
Berapa menit: setengah jam; \omterm{def:g3:division:half}{seperempat} jam; tiga perempat jam;
$\frac13$ jam?
\end{exercise}

\begin{exercise}[$\star$]\label{exo:g4:fractions:9}
Di sebuah kelas berisi $24$ murid, $\frac14$ memakai kacamata.
Berapa murid yang memakai kacamata? Berapa yang tidak?
\end{exercise}

\begin{exercise}[$\star\star$]\label{exo:g4:fractions:10}
Nina memakan $\frac38$ bagian sebuah pizza dan Sandi memakan
$\frac48$ bagian pizza yang sama. Berapa bagian yang mereka makan
bersama? Berapa bagian yang tersisa? Siapa yang makan lebih banyak?
\end{exercise}

\begin{exercise}[$\star\star$]\label{exo:g4:fractions:11}
Benar atau salah, dengan gambar atau contoh: ``mungkin saja
$\frac12$ pizza kecil lebih sedikit daripada $\frac14$ pizza yang
sangat besar''. Apa yang harus selalu diketahui sebelum
membandingkan dua \omterm{def:g4:fractions:def}{pecahan} dari sesuatu?

\end{exercise}
